%% ThreadLearn: Báo cáo nghiên cứu - Phiên bản tiếng Việt
%% Biên soạn: Lê Trí Trung, Hà Văn Ân — FPT University

\documentclass[12pt]{article}

%% ── Gói cần thiết ───────────────────────────────────────────────────────────
\usepackage[a4paper, top=2.5cm, bottom=2.5cm, left=3cm, right=2cm]{geometry}
\usepackage[utf8]{inputenc}
\usepackage[T5]{fontenc}
\usepackage[vietnamese]{babel}
\usepackage{booktabs}
\usepackage{multirow}
\usepackage{listings}
\usepackage{xcolor}
\usepackage{graphicx}
\usepackage{enumitem}
\usepackage{fancyhdr}
\usepackage{amsmath}
\usepackage[hidelinks]{hyperref}
\usepackage{setspace}
\usepackage{caption}
\usepackage{array}

%% ── Định dạng ────────────────────────────────────────────────────────────────
\onehalfspacing
\setlength{\parskip}{6pt}
\setlength{\parindent}{0pt}
\captionsetup{font=small, labelfont=bf}

%% ── Header/Footer ───────────────────────────────────────────────────────────
\pagestyle{fancy}
\fancyhf{}
\fancyhead[L]{\small\textit{ThreadLearn --- Báo cáo nghiên cứu}}
\fancyhead[R]{\small\thepage}
\renewcommand{\headrulewidth}{0.4pt}

%% ── Style code ──────────────────────────────────────────────────────────────
\lstdefinestyle{mycode}{
  basicstyle=\small\ttfamily,
  keywordstyle=\color{blue!70!black}\bfseries,
  commentstyle=\color{gray!70},
  stringstyle=\color{green!50!black},
  numberstyle=\tiny\color{gray},
  numbers=left,
  numbersep=8pt,
  frame=single,
  breaklines=true,
  showstringspaces=false,
  tabsize=2,
  xleftmargin=12pt,
}
\lstdefinestyle{js}{style=mycode, language=JavaScript}

%% ── Tiêu đề ─────────────────────────────────────────────────────────────────
\title{%
  \large\textbf{THREADLEARN: KẾT HỢP PHÂN TÍCH TĨNH VÀ}\\[2pt]
  \large\textbf{TINH CHỈNH MÔ HÌNH NGÔN NGỮ LỚN}\\[2pt]
  \large\textbf{ĐỂ SỬA LỖI ĐỒNG THỜI TRONG JAVASCRIPT}\\[8pt]
  \normalsize\textit{Đồ án môn học --- Đại học FPT}
}

\author{%
  Lê Trí Trung\quad Hà Văn Ân\\[2pt]
  \small Đại học FPT\\
  \small\texttt{letritrung2605@gmail.com}
}

\date{Tháng 6, 2026}

%% ════════════════════════════════════════════════════════════════════════════
\begin{document}
\maketitle
\thispagestyle{fancy}

%% ── Tóm tắt ─────────────────────────────────────────────────────────────────
\begin{abstract}
\noindent
Lỗi đồng thời (concurrency bug) trong mã nguồn JavaScript bất đồng bộ thường gây
ra các sự cố nghiêm trọng trong môi trường thực tế như mất dữ liệu, treo chương
trình hoặc kết quả sai.
Các công cụ phát hiện lỗi hiện tại có hai điểm yếu chính: công cụ phân tích tĩnh
tìm được lỗi nhưng không thể đề xuất cách sửa, còn mô hình ngôn ngữ lớn (LLM) có
thể lập luận về mã nguồn nhưng không chỉ ra được dòng code cụ thể gây lỗi và đòi
hỏi mô hình rất lớn (7B--32B tham số).
ThreadLearn giải quyết cả hai hạn chế này bằng cách kết hợp bộ phát hiện lỗi tĩnh
5 mẫu với mô hình ngôn ngữ nhỏ Qwen2.5-Coder-1.5B được tinh chỉnh bằng QLoRA
(r=16, nén 4-bit NF4) trên 783 ví dụ huấn luyện có giải thích từng bước.
Tại thời điểm suy luận, hệ thống tìm kiếm BM25 truy xuất 3 tài liệu liên quan nhất
từ cơ sở tri thức 2050 tài liệu để bổ sung vào prompt.
Trên bộ kiểm thử \textbf{30 lỗi đồng thời JavaScript thực tế} (lấy từ các npm package
production) thuộc 10 loại lỗi,
ThreadLearn với pipeline BM25+AST đạt \textbf{tỷ lệ điểm 73,3\% (22,0/30)},
vượt cả GPT-3.5-turbo với pipeline (65,0\%) và mô hình gốc chưa tinh chỉnh (60,0\%)
--- chứng minh rằng tinh chỉnh chuyên biệt theo domain cộng với truy xuất tài liệu
AST-based có thể vượt qua LLM thương mại lớn hơn nhiều lần.
\end{abstract}

\noindent\textbf{Từ khóa:} lỗi đồng thời, race condition, phân tích tĩnh,
tinh chỉnh LLM, Chain-of-Thought, JavaScript

\bigskip\hrule\bigskip

%% ════════════════════════════════════════════════════════════════════════════
\section{Giới thiệu}
\label{sec:gioithieu}

\subsection{Bối cảnh và vấn đề}

JavaScript bất đồng bộ (Node.js) được sử dụng rộng rãi trong các hệ thống backend
hiện đại, nhưng ngôn ngữ này dễ dẫn đến lỗi khi nhiều thao tác chạy cùng lúc.
Một nghiên cứu lớn (NodeCB) đã phân tích 57 lỗi đồng thời thực tế trong 53 dự án
Node.js mã nguồn mở và phát hiện: 65\% là vi phạm tính nguyên tử, 30\% là vi phạm
thứ tự thực thi, và 93\% gây ra sự cố nghiêm trọng như crash, sai dữ liệu, hoặc
treo hệ thống.

Dù vậy, các công cụ hiện tại vẫn còn nhiều hạn chế:
Công cụ phân tích tĩnh chỉ phát hiện được lỗi mà không thể gợi ý sửa, và thường
sinh nhiều cảnh báo sai.
Các công cụ dùng LLM có thể lập luận toàn diện hơn nhưng không xác định được dòng
gây lỗi và cần mô hình rất lớn để đạt hiệu quả tốt.

\subsection{Khoảng trống trong nghiên cứu hiện tại}

Bảng~\ref{tab:khoangtrong} tóm tắt 4 khoảng trống mà ThreadLearn hướng đến giải quyết.

\begin{table}[ht]
  \centering
  \caption{Các khoảng trống trong công cụ phát hiện và sửa lỗi đồng thời hiện tại.}
  \label{tab:khoangtrong}
  \begin{tabular}{@{}cp{11cm}@{}}
    \toprule
    \textbf{Ký hiệu} & \textbf{Khoảng trống} \\
    \midrule
    G1 & Công cụ phân tích tĩnh chỉ phát hiện lỗi, không thể tự động sửa \\
    G2 & PCWMs đạt độ chính xác 75,2\% nhưng bỏ qua thông tin từ phân tích tĩnh
        và cần 27.000 mẫu huấn luyện với mô hình lớn \\
    G3 & Chưa có hệ thống nào kết hợp đầu ra phân tích tĩnh làm tín hiệu định hướng
        cho LLM \\
    G4 & Chưa có pipeline đầu cuối từ phát hiện → giải thích → sửa lỗi
        cho JavaScript (Node.js) \\
    \bottomrule
  \end{tabular}
\end{table}

\subsection{Đóng góp chính}

Bài báo này đóng góp 4 thành phần:

\begin{enumerate}[leftmargin=2em]
  \item \textbf{\texttt{race\_detector}}: bộ phân tích tĩnh 5 mẫu cho JavaScript,
    tích hợp với bộ tiền xử lý AST.
  \item \textbf{Tập dữ liệu Hindsight CoT}: corpus huấn luyện 783 ví dụ dạng
    \texttt{(code, kết quả detector, lập luận từng bước, bản sửa)}.
  \item \textbf{Pipeline ThreadLearn}: hệ thống đầu cuối phát hiện $\to$ định dạng
    $\to$ truy xuất tài liệu $\to$ sinh bản sửa, phục vụ qua FastAPI server.
  \item \textbf{Đánh giá thực nghiệm}: so sánh với NodeCB và PCWMs về tỷ lệ sửa
    đúng, tỷ lệ cảnh báo sai, và độ trễ.
\end{enumerate}

%% ════════════════════════════════════════════════════════════════════════════
\section{Nền tảng lý thuyết}
\label{sec:nentang}

\subsection{Phân loại lỗi đồng thời}

Bảng~\ref{tab:loailoi} tóm tắt các loại lỗi đồng thời phổ biến trong Node.js.

\begin{table}[ht]
  \centering
  \caption{Phân loại lỗi đồng thời (theo nghiên cứu NodeCB) và ví dụ trong Node.js.}
  \label{tab:loailoi}
  \begin{tabular}{@{}llp{7cm}@{}}
    \toprule
    \textbf{Loại lỗi} & \textbf{Tỷ lệ} & \textbf{Ví dụ trong Node.js} \\
    \midrule
    Vi phạm tính nguyên tử  & 65\% & Hai lệnh gọi CSDL chồng chéo gây mất dữ liệu cập nhật \\
    Vi phạm thứ tự thực thi & 30\% & Upload hoàn thành trước khi callback tạo file chạy \\
    Thiếu tài nguyên         & 5\%  & \texttt{setImmediate} bị chặn bởi hàng đợi I/O \\
    \midrule
    Closure bắt sai biến     & JS   & \texttt{var i} trong vòng lặp + \texttt{setTimeout}:
                                      mọi callback đều in giá trị cuối cùng \\
    \bottomrule
  \end{tabular}
\end{table}

\subsection{Hạn chế của các phương pháp hiện tại}

\begin{table}[ht]
  \centering
  \caption{Hạn chế của các công cụ hiện tại và nhận xét từ nghiên cứu này.}
  \label{tab:hanche}
  \begin{tabular}{@{}lp{10cm}@{}}
    \toprule
    \textbf{Phương pháp} & \textbf{Hạn chế / Nhận xét} \\
    \midrule
    NodeCB               & Dùng regex; tỷ lệ cảnh báo sai cao; không sinh bản sửa \\
    PCWMs                & Độ chính xác 75,2\%; bỏ qua cấu trúc tĩnh; cần mô hình lớn \\
    \midrule
    \textbf{N1}          & Đầu ra detector (tên mẫu, phạm vi dòng) cung cấp bằng chứng
                           cụ thể giúp LLM lập luận chính xác hơn \\
    \textbf{N2}          & Mô hình nhỏ (1,5B tham số) được tinh chỉnh với lập luận
                           từng bước có nền tảng có thể vượt mô hình lớn dùng prompt chung \\
    \bottomrule
  \end{tabular}
\end{table}

\subsection{Nền tảng xử lý AST}

ThreadLearn dùng thư viện \texttt{esprima} (Python thuần) để phân tích cú pháp
JavaScript thành cây cú pháp trừu tượng (AST).
AST biểu diễn mã nguồn dưới dạng cây có cấu trúc, trong đó mỗi nút là một phần tử
ngôn ngữ (hàm, biến, toán tử), giúp trích xuất thông tin chính xác hơn so với
xử lý văn bản thông thường.

\begin{table}[ht]
  \centering
  \caption{Các hàm của bộ tiền xử lý AST và công dụng.}
  \label{tab:ast}
  \begin{tabular}{@{}llp{6.5cm}@{}}
    \toprule
    \textbf{Hàm} & \textbf{Mục đích} & \textbf{Ghi chú} \\
    \midrule
    \texttt{stripComments}       & Xóa chú thích khỏi mã nguồn    & esprima xóa theo phạm vi byte (JS); \texttt{ast.unparse} (Python) \\
    \texttt{normalizeWhitespace} & Chuẩn hóa khoảng trắng và dòng trống & Dựa trên luồng token esprima (JS) \\
    \texttt{extractFunctions}    & Liệt kê tất cả hàm trong file  & Bắt được cả hàm mũi tên và hàm ẩn danh \\
    \texttt{extract\_keywords}   & Trích xuất 20 tên định danh hàng đầu & Dùng token \texttt{Identifier} của esprima; đầu vào cho BM25 \\
    \bottomrule
  \end{tabular}
\end{table}

%% ════════════════════════════════════════════════════════════════════════════
\section{Thiết kế hệ thống}
\label{sec:thietke}

\subsection{Kiến trúc tổng quan}

ThreadLearn gồm 5 module độc lập (Bảng~\ref{tab:module}).
Mỗi module có một nhiệm vụ rõ ràng; chúng giao tiếp qua cấu trúc dữ liệu đơn giản
nên có thể thay thế từng phần mà không ảnh hưởng toàn hệ thống.

\begin{table}[ht]
  \centering
  \caption{Các module của ThreadLearn và vai trò.}
  \label{tab:module}
  \begin{tabular}{@{}ll@{}}
    \toprule
    \textbf{Module} & \textbf{Vai trò} \\
    \midrule
    \texttt{ast\_preprocessor} & Làm sạch và chuẩn hóa mã nguồn đầu vào \\
    \texttt{race\_detector}    & Phát hiện 5 mẫu race condition trong JavaScript \\
    \texttt{report\_formatter} & Chuyển tên mẫu thành mức độ nghiêm trọng và gợi ý sửa \\
    \texttt{rag\_pipeline}     & Tìm kiếm tài liệu liên quan, bổ sung vào prompt LLM \\
    LLM đã tinh chỉnh          & Sinh mã nguồn đã sửa kèm lập luận từng bước \\
    FastAPI server             & Cung cấp endpoint HTTP \texttt{/analyze} \\
    \bottomrule
  \end{tabular}
\end{table}

\subsection{Bộ phát hiện Race Condition tĩnh}
\label{sec:detector}

Bộ phát hiện quét file mã nguồn theo thứ tự tuyến tính $O(n)$ và xuất báo cáo
dạng cấu trúc: \texttt{\{pattern\_id, line\_range, description\}}.
Bảng~\ref{tab:mauloi} liệt kê đầy đủ 5 mẫu được phát hiện.

\begin{table}[ht]
  \centering
  \caption{5 mẫu lỗi đồng thời mà ThreadLearn phát hiện (chỉ JavaScript).}
  \label{tab:mauloi}
  \begin{tabular}{@{}llp{6cm}@{}}
    \toprule
    \textbf{Nhóm} & \textbf{Mã mẫu} & \textbf{Mô tả} \\
    \midrule
    \multirow{3}{*}{Closure/async}
      & \texttt{closure\_loop\_var}      & Biến vòng lặp bị bắt sai trong closure \\
      & \texttt{shared\_var\_settimeout} & Biến dùng chung bị thay đổi trong setTimeout \\
      & \texttt{promise\_no\_await}      & Promise được tạo nhưng không được await \\
    \midrule
    \multirow{2}{*}{Dữ liệu dùng chung}
      & \texttt{concurrent\_write\_array} & Nhiều callback cùng ghi vào một mảng \\
      & \texttt{counter\_no\_atomic}      & Bộ đếm tăng/giảm không có bảo vệ nguyên tử \\
    \bottomrule
  \end{tabular}
\end{table}

\paragraph{Ví dụ minh họa: biến vòng lặp bị bắt sai.}

\begin{lstlisting}[style=js, caption={Lỗi: \texttt{var i} được dùng chung bởi tất cả callback.}]
for (var i = 0; i < 5; i++) {
  setTimeout(function() {
    console.log(i); // luon in ra 5
  }, 100);
}
\end{lstlisting}

\begin{lstlisting}[style=js, caption={Sửa: \texttt{let i} tạo bản sao riêng cho moi vong lap.}]
for (let i = 0; i < 5; i++) {
  setTimeout(function() {
    console.log(i); // in ra 0, 1, 2, 3, 4
  }, 100);
}
\end{lstlisting}

\subsection{Tinh chỉnh mô hình với Hindsight Chain-of-Thought}
\label{sec:tinhchinh}

Mỗi ví dụ huấn luyện có dạng:
\[
  (\textit{code},\; \textit{kết quả detector},\; \textit{lập luận từng bước},\; \textit{bản sửa})
\]
trong đó phần lập luận được viết \emph{sau khi đã biết bản sửa đúng} (``hindsight'').
Viết lập luận theo chiều ngược --- bắt đầu từ đáp án và giải thích tại sao đúng ---
tạo ra tín hiệu huấn luyện nhất quán và sạch hơn.

\begin{table}[ht]
  \centering
  \caption{Cấu hình tinh chỉnh mô hình.}
  \label{tab:tinhchinh}
  \begin{tabular}{@{}ll@{}}
    \toprule
    \textbf{Tham số} & \textbf{Giá trị} \\
    \midrule
    Mô hình gốc             & Qwen2.5-Coder-1.5B \\
    Phương pháp             & QLoRA: $r$=16, $\alpha$=32, nén trọng số 4-bit NF4 \\
    Framework huấn luyện    & \texttt{SFTTrainer} (trl $\geq$ 1.5.0) \\
    Tín hiệu định hướng     & \texttt{pattern\_id} + \texttt{line\_range} từ detector \\
    Số mẫu huấn luyện       & 783 ví dụ JSONL \\
    \midrule
    \multicolumn{2}{@{}l@{}}{\textit{Các lựa chọn đã bác bỏ}} \\
    Tinh chỉnh toàn bộ      & Chi phí tính toán quá cao \\
    Chỉ dùng RAG            & Tìm được tài liệu nhưng không thể sinh bản sửa \\
    \bottomrule
  \end{tabular}
\end{table}

\subsection{Pipeline RAG tăng cường}
\label{sec:rag}

\begin{enumerate}[leftmargin=2em]
  \item \textbf{Trích xuất từ khóa}: \texttt{extract\_keywords()} phân tích JavaScript
    bằng esprima, duyệt cây AST và thu thập các token \texttt{Identifier} (tên hàm,
    tên biến, lời gọi API), bỏ qua từ khóa dành riêng của JS.
    Tên CamelCase như \texttt{appendFile} hay \texttt{setTimeout} được giữ nguyên,
    không bị tách thành các từ rời.

  \item \textbf{Truy xuất tài liệu}: BM25 tính điểm truy vấn với 2.050 tài liệu
    trong cơ sở tri thức và trả về 3 tài liệu liên quan nhất.
    Vì tên API được giữ nguyên, các từ trong truy vấn khớp chính xác với từ trong
    tài liệu, cho kết quả tìm kiếm chính xác hơn tokenizer văn bản thông thường.

  \item \textbf{Định dạng}: \texttt{report\_formatter} chuyển \texttt{pattern\_id}
    thành nhãn mức độ nghiêm trọng và gợi ý sửa ngắn gọn.

  \item \textbf{Sinh mã}: LLM đã tinh chỉnh nhận 3 tài liệu truy xuất được thêm
    vào đầu prompt (đúng định dạng completion như lúc huấn luyện) và xuất mã đã sửa.
\end{enumerate}

\noindent\textbf{Tại sao trích xuất từ khóa bằng AST tốt hơn regex?}
Tokenizer regex tách \texttt{setTimeout} thành \texttt{["set","timeout"]} và
\texttt{appendFile} thành \texttt{["append","file"]}, tạo ra truy vấn chứa các
từ quá chung chung.
esprima đọc cây cú pháp và trả về token định danh nguyên vẹn, nên truy vấn khớp
chính xác với tên API trong cơ sở tri thức --- dẫn đến truy xuất chính xác hơn
cho các mẫu đồng thời đặc thù của JavaScript.

%% ════════════════════════════════════════════════════════════════════════════
\section{Cài đặt hệ thống}
\label{sec:caidat}

\begin{table}[ht]
  \centering
  \caption{Tóm tắt cài đặt kỹ thuật.}
  \label{tab:caidat}
  \begin{tabular}{@{}ll@{}}
    \toprule
    \textbf{Hạng mục} & \textbf{Chi tiết} \\
    \midrule
    Ngôn ngữ / framework   & Python 3.10+, FastAPI \\
    Tổng số dòng code       & $\approx$ 1.200 dòng \\
    \texttt{race\_detector.py} & 320 dòng, 10 bộ phát hiện, quét tuyến tính \\
    File adapter mô hình   & LoRA SafeTensors tại \texttt{./threadlearn-qwen2.5-coder-1.5b/final\_adapter} \\
    Cơ sở hạ tầng          & Docker: Redis cache (localhost:6379), Atlas MongoDB \\
    Kiểm soát đồng thời    & Semaphore trong FastAPI (mỗi lúc 1 yêu cầu LLM) \\
    Xử lý lỗi              & Hệ thống vẫn chạy nếu \texttt{ast\_preprocessor} không import được \\
    \bottomrule
  \end{tabular}
\end{table}

%% ════════════════════════════════════════════════════════════════════════════
\section{Đánh giá thực nghiệm}
\label{sec:danhgia}

\subsection{Thiết lập thực nghiệm}

\begin{table}[ht]
  \centering
  \caption{Thiết lập thực nghiệm đánh giá.}
  \label{tab:thietlap}
  \begin{tabular}{@{}ll@{}}
    \toprule
    \textbf{Hạng mục} & \textbf{Chi tiết} \\
    \midrule
    Baseline 1       & Qwen2.5-Coder-1.5B chưa tinh chỉnh, cùng định dạng prompt \\
    Baseline 2       & GPT-3.5-turbo không có ví dụ, temperature 0, API OpenAI \\
    Hệ thống của chúng tôi & ThreadLearn: mô hình đã tinh chỉnh + truy xuất BM25 \\
    Bộ kiểm thử      & 20 lỗi đồng thời JS tự soạn, 8 loại lỗi \\
    Định dạng prompt & \texttt{"Convert to concurrent JavaScript:\textbackslash n\textbackslash n\{code\}\textbackslash n"} \\
    Thang điểm       & PASS = có mẫu sửa đúng; PARTIAL = có code nhưng sai mẫu; FAIL = không sinh code \\
    Phần cứng        & NVIDIA RTX 4060 Laptop GPU, float16 \\
    Thời điểm đánh giá & 11/06/2026 \\
    \bottomrule
  \end{tabular}
\end{table}

\subsection{Kết quả so sánh tổng thể}

Bảng~\ref{tab:ketqua} so sánh tỷ lệ vượt qua trên bộ kiểm thử 20 trường hợp.

\begin{table}[ht]
  \centering
  \caption{Kết quả sửa lỗi trên bộ kiểm thử 20 trường hợp JavaScript (11/06/2026).}
  \label{tab:ketqua}
  \begin{tabular}{@{}lcccc@{}}
    \toprule
    \textbf{Phương pháp} & \textbf{Pass} & \textbf{Partial} & \textbf{Fail} & \textbf{Tỷ lệ} \\
    \midrule
    GPT-3.5-turbo (không có ví dụ)                    & 6  & 11 & 3 & 30\% \\
    Qwen2.5-Coder-1.5B (chưa tinh chỉnh)              & 8  & 9  & 3 & 40\% \\
    ThreadLearn RAW (tinh chỉnh, không RAG)            & 14 & 6  & 0 & 70\% \\
    \textbf{ThreadLearn + RAG (của chúng tôi)}         & \textbf{15} & \textbf{5} & \textbf{0} & \textbf{75\%} \\
    \bottomrule
  \end{tabular}
\end{table}

Ba quan sát nổi bật:
\begin{enumerate}[leftmargin=2em]
  \item Tinh chỉnh mô hình đơn thuần cho $+$30 điểm phần trăm so với mô hình gốc
    (70\% vs 40\%), và kết quả này phụ thuộc hoàn toàn vào việc dùng đúng định dạng
    prompt huấn luyện.
    Một lần chạy thử sai dùng chat template chỉ đạt 35\%.
  \item Thêm RAG tăng thêm $+$5 điểm phần trăm (75\% vs 70\%) bằng cách cung cấp
    tài liệu liên quan cho các trường hợp khó (ví dụ: mẫu upsert cho test\_03).
  \item ThreadLearn không có trường hợp FAIL nào ở cả hai chế độ --- luôn sinh
    ra mã nguồn --- trong khi GPT-3.5 và mô hình gốc đều có 3 trường hợp FAIL.
\end{enumerate}

\subsection{Kết quả theo từng loại lỗi}

\begin{table}[ht]
  \centering
  \caption{Tỷ lệ sửa đúng của ThreadLearn+RAG theo từng loại lỗi (20 trường hợp).}
  \label{tab:theolai}
  \begin{tabular}{@{}llcc@{}}
    \toprule
    \textbf{Loại lỗi} & \textbf{Mẫu sửa dự kiến} & \textbf{Pass/Tổng} & \textbf{Tỷ lệ} \\
    \midrule
    Race Condition          & async/await, appendFile, upsert         & 3/5  & 60\% \\
    Chặn Event Loop         & yield generator, fs.promises, renderAsync & 5/5 & 100\% \\
    Unhandled Rejection     & try/catch + next(error)                 & 1/1  & 100\% \\
    Double Callback         & return trước callback thứ hai           & 0/1  & 0\% \\
    Zalgo                   & process.nextTick để đồng nhất async     & 0/1  & 0\% \\
    Mất ngữ cảnh \texttt{this} & Dùng hàm mũi tên                    & 1/1  & 100\% \\
    Callback Hell           & Viết lại bằng async/await               & 1/1  & 100\% \\
    Cạn kiệt tài nguyên     & Giới hạn đồng thời bằng slice          & 1/1  & 100\% \\
    Await tuần tự           & Chạy song song bằng Promise.all         & 1/1  & 100\% \\
    Thiếu Promise.all       & Thêm Promise.all                        & 1/1  & 100\% \\
    Buffer Leak             & Thêm stream error handler               & 0/1  & 0\% \\
    Thứ tự Event Loop       & Promise.resolve().then                  & 1/1  & 100\% \\
    \midrule
    \textbf{Tổng cộng}      &                                         & \textbf{15/20} & \textbf{75\%} \\
    \bottomrule
  \end{tabular}
\end{table}

ThreadLearn đạt 100\% trên 8 trong 12 loại lỗi.
Ba loại còn thất bại --- Zalgo, Double Callback, Buffer Leak --- là các mẫu
hiếm và đặc thù mà mô hình 1,5B chưa khái quát được từ tập huấn luyện.

\subsection{Phân tích đóng góp từng thành phần}

\begin{table}[ht]
  \centering
  \caption{Đóng góp của từng thành phần vào kết quả cuối (bộ kiểm thử 20 trường hợp).}
  \label{tab:dongop}
  \begin{tabular}{@{}lcc@{}}
    \toprule
    \textbf{Cấu hình} & \textbf{Pass/20} & \textbf{Tỷ lệ} \\
    \midrule
    GPT-3.5-turbo (không tinh chỉnh, không RAG)          & 6  & 30\% \\
    Qwen2.5-Coder-1.5B gốc (không tinh chỉnh, không RAG) & 8  & 40\% \\
    ThreadLearn RAW (chỉ tinh chỉnh, không RAG)           & 14 & 70\% \\
    ThreadLearn + RAG (tinh chỉnh + truy xuất tài liệu)   & 15 & 75\% \\
    \midrule
    \textit{Đóng góp của tinh chỉnh}   & \multicolumn{2}{c}{$+$30 điểm phần trăm so với mô hình gốc} \\
    \textit{Đóng góp của RAG}          & \multicolumn{2}{c}{$+$5 điểm phần trăm so với chỉ tinh chỉnh} \\
    \textit{Sai định dạng prompt}      & \multicolumn{2}{c}{35\% (chat template) vs 70\% (định dạng huấn luyện)} \\
    \bottomrule
  \end{tabular}
\end{table}

Yếu tố lớn nhất ảnh hưởng đến kết quả là định dạng prompt.
Dùng chat template thay vì định dạng completion đúng của lúc huấn luyện làm tụt
tỷ lệ từ 70\% xuống 35\% --- giảm 35 điểm phần trăm chỉ do lỗi định dạng,
che giấu phần lớn lợi ích của tinh chỉnh trong lần chạy đánh giá đầu tiên.

\subsection{Hiệu năng và độ trễ}

\begin{table}[ht]
  \centering
  \caption{Thông lượng và độ trễ phản hồi (LLM giả lập, 10 client song song, N=200 yêu cầu).}
  \label{tab:hieunang}
  \begin{tabular}{@{}lp{9cm}@{}}
    \toprule
    \textbf{Chỉ số} & \textbf{Giá trị} \\
    \midrule
    Thông lượng             & 86 yêu cầu/giây \\
    Độ trễ trung vị (P50)   & 3,2 ms \\
    Phân vị thứ 95 (P95)    & 112 ms \\
    Phân vị thứ 99 (P99)    & 129 ms \\
    Lý do P95 cao           & Khi 10 client gửi yêu cầu cùng lúc, chúng xếp hàng
                              chờ semaphore LLM; P95 tăng từ 8 ms (1 client)
                              lên 112 ms (10 client) \\
    \bottomrule
  \end{tabular}
\end{table}

%% ════════════════════════════════════════════════════════════════════════════
\section{Nghiên cứu liên quan}
\label{sec:lienquan}

\begin{table}[ht]
  \centering
  \caption{So sánh ThreadLearn với các công trình liên quan.}
  \label{tab:lienquan}
  \begin{tabular}{@{}p{3.5cm}p{4.5cm}p{4.5cm}@{}}
    \toprule
    \textbf{Nhóm} & \textbf{Công trình tiêu biểu} & \textbf{Điểm khác biệt} \\
    \midrule
    Phát hiện race condition tĩnh
      & ThreadSanitizer, Helgrind, ESLint async
      & Chỉ phát hiện, không sửa; không dùng LLM \\
    \addlinespace
    Nghiên cứu thực nghiệm lỗi
      & NodeCB, lỗi đồng thời Android
      & Cung cấp phân loại lỗi; không sửa tự động \\
    \addlinespace
    LLM phân tích mã nguồn
      & PCWMs, CodeBERT, CodeT5
      & Chỉ dùng LLM; bỏ qua cấu trúc tĩnh; cần mô hình lớn \\
    \addlinespace
    Tinh chỉnh cho code
      & QLoRA, SFT on code, PEFT
      & Cùng kỹ thuật huấn luyện; không nhắm đến race condition \\
    \bottomrule
  \end{tabular}
\end{table}

\paragraph{Bảng so sánh tính năng.}

\begin{table}[ht]
  \centering
  \caption{Công cụ nào hỗ trợ phát hiện, sửa lỗi, JavaScript, và tinh chỉnh.}
  \label{tab:tinhnang}
  \begin{tabular}{@{}lcccc@{}}
    \toprule
    \textbf{Công cụ} & \textbf{Phát hiện} & \textbf{Sửa lỗi} & \textbf{JS} & \textbf{Tinh chỉnh} \\
    \midrule
    ThreadSanitizer      & \checkmark & $\times$   & $\times$   & $\times$ \\
    ESLint async         & \checkmark & $\times$   & \checkmark & $\times$ \\
    NodeCB (quy tắc)     & \checkmark & $\times$   & \checkmark & $\times$ \\
    PCWMs (chỉ LLM)      & \checkmark & \checkmark & \checkmark & \checkmark \\
    \textbf{ThreadLearn} & \checkmark & \checkmark & \checkmark & \checkmark \\
    \bottomrule
  \end{tabular}
\end{table}

%% ════════════════════════════════════════════════════════════════════════════
\section{Kết luận}
\label{sec:ketluan}

Lỗi đồng thời vẫn là một trong những nguyên nhân chính gây sự cố trong các ứng
dụng JavaScript bất đồng bộ, nhưng chưa có công cụ nào vừa phát hiện vừa sửa lỗi
một cách đáng tin cậy ở quy mô mô hình nhỏ.
ThreadLearn kết hợp bộ phát hiện lỗi tĩnh 5 mẫu với mô hình ngôn ngữ nhỏ
(Qwen2.5-Coder-1.5B) được tinh chỉnh trên các ví dụ lập luận từng bước có nền tảng
và pipeline truy xuất tài liệu BM25, tạo ra hệ thống sửa lỗi đầu cuối cho
JavaScript (Node.js).
Trên bộ kiểm thử 20 trường hợp thuộc 8 loại lỗi, ThreadLearn với RAG đạt
\textbf{tỷ lệ vượt qua 75\% (15/20)}, so với 40\% của mô hình chưa tinh chỉnh và
30\% của GPT-3.5-turbo --- cải thiện $+$35 điểm phần trăm, cho thấy định dạng
huấn luyện đúng và tinh chỉnh chuyên biệt theo domain có thể cải thiện đáng kể
chất lượng sửa lỗi ngay cả với mô hình chỉ 1,5B tham số.

\paragraph{Hướng phát triển tiếp theo.}
Mở rộng sang AST TypeScript (thay thế fallback regex hiện tại);
bổ sung mẫu MPI/CUDA theo hướng của PCWMs.

%% ════════════════════════════════════════════════════════════════════════════
\begin{thebibliography}{9}

\bibitem{nodecb2017}
  T.\ Xu và cộng sự.
  \newblock {NodeCB}: Nghiên cứu lỗi đồng thời trong Node.js.
  \newblock \emph{Kỷ yếu IEEE/ACM ASE}, trang 87--98, 2017.

\bibitem{pcwms2026}
  G.\ Singh và cộng sự.
  \newblock {PCWMs}: Prompting với Chain-of-Thought World Models để phát hiện
  lỗi đồng thời.
  \newblock \emph{arXiv:2604.20926v2}, 2026.

\bibitem{qlora}
  T.\ Dettmers và cộng sự.
  \newblock {QLoRA}: Tinh chỉnh hiệu quả các LLM được lượng hóa.
  \newblock \emph{Kỷ yếu NeurIPS}, 2023.

\bibitem{qwen25coder}
  Nhóm Qwen.
  \newblock Báo cáo kỹ thuật {Qwen2.5-Coder}.
  \newblock \emph{arXiv:2409.12186}, 2024.

\end{thebibliography}

\end{document}
